\chapter{How to Use This Book}
\label{ch:how-to-use}

\section*{Five parts, one dependency order}

The book has five parts. Part~I is the foundation and everything else assumes it.
Parts~II and~III are the core; they can be read in either order after Part~I.
Part~IV applies the model operationally, and Part~V is optional depth.
\Cref{fig:how-parts} shows the dependencies.

\begin{figure}[htbp]
\centering
\begin{tikzpicture}[
  node distance=0.9cm and 1.1cm,
  part/.style={draw=#1, fill=#1!12, rounded corners=3pt, line width=1pt,
               minimum width=3.4cm, minimum height=1.25cm, align=center,
               font=\sffamily\small, text=bkGrayInk},
  arr/.style={-{Stealth[length=2.6mm]}, line width=1pt, bkGrayInk!70}]
  \node[part=bkBlue]   (p1) {\textbf{Part I}\\How Spark executes\\a query};
  \node[part=bkGreen, above right=0.45cm and 1.2cm of p1] (p2) {\textbf{Part II}\\Making queries\\fast};
  \node[part=bkPurple, below right=0.45cm and 1.2cm of p1] (p3) {\textbf{Part III}\\The Iceberg\\lakehouse};
  \node[part=bkAmber,  right=1.2cm of $(p2)!0.5!(p3)$] (p4) {\textbf{Part IV}\\Operating\\Spark};
  \node[part=bkRose,   right=1.0cm of p4] (p5) {\textbf{Part V}\\Internals\\(optional)};
  \draw[arr] (p1) -- (p2);
  \draw[arr] (p1) -- (p3);
  \draw[arr] (p2) -- (p4);
  \draw[arr] (p3) -- (p4);
  \draw[arr, dashed] (p4) -- (p5);
\end{tikzpicture}
\caption{Reading order. Part~I first; Parts~II and~III in either order; Part~IV
once you have the model; Part~V when you want to know how it is built.}
\label{fig:how-parts}
\end{figure}

Within a part, chapters build on their predecessors, but each states its
prerequisites at the start so you can jump in.

\section*{Four ways to read it}

\begin{description}[leftmargin=0pt, labelindent=0pt, itemsep=6pt]
  \item[Straight through.] The intended path. Each chapter ends with a summary and
  exercises, and the running specimen carries from one to the next.
  \item[``My query is slow right now.''] Read \Cref{ch:shape} for the vocabulary,
  then the \texttt{EXPLAIN} sections of \Cref{ch:catalyst} and \Cref{app:plans} to
  read the plan, then \Cref{ch:debugging} to localize the cost. Each chapter's
  symptom table (see below) then points you to the mechanism that explains it.
  \item[``I run Iceberg and want the lakehouse side first.''] \Cref{ch:shape} and
  \Cref{ch:pruning} for the Spark-side vocabulary of pruning, then Part~III in
  order. Return to Part~II when a chapter in Part~III refers back to it.
  \item[``I size and operate clusters.''] \Cref{ch:execmem}, \Cref{ch:shuffle}, and
  Part~IV. Skip the join and Iceberg chapters until a symptom sends you there.
\end{description}

\section*{The shape of a chapter}

The chapters share one pattern, so you know where to look:

\begin{enumerate}
  \item an \textbf{opening} that states the question the chapter answers, using the
  running example;
  \item the \textbf{mechanism}, explained before any setting is mentioned;
  \item a \textbf{measured demonstration} --- the plan, the metric, or the file
  listing --- usually from the lab (\Cref{sec:refenv-lab});
  \item a \textbf{symptom $\rightarrow$ evidence $\rightarrow$ knob $\rightarrow$ cost}
  table, which is the chapter's field guide: what you will observe, where to see
  it, what to change, and what the change costs;
  \item a \textbf{summary and exercises}.
\end{enumerate}

The tables are deliberately last. If you start with them, you get the configs
without the model, which is the failure mode this book is written against.

\section*{The running example}

One dataset runs through every chapter: an online marketplace's \texttt{orders}
table, an Iceberg table partitioned by day, with a small star schema around it
(\texttt{dim\_product}, \texttt{dim\_customer}, \texttt{dim\_date}). It is
defined in \Cref{app:refenv}. Using the same tables everywhere means a change in
one chapter can be seen in the plans of another, and you never have to learn a new
toy schema.

\section*{Running the examples yourself}

The book does not require a cluster. The \texttt{lab/} directory of the
repository contains a Docker-based environment with Spark 3.5.2, the Iceberg
runtime, a generator for the synthetic \texttt{orders} data, and one script per
experiment; \Cref{sec:refenv-lab} describes the setup and its limits. The lab is
a single machine, so it reproduces the \emph{mechanisms} and the \emph{ratios}
reliably and the absolute timings not at all. Where a chapter reports a time, it
says which environment produced it.

If your stack differs from the reference --- another Spark version, a different
catalog or object store --- the mechanisms in the
book still hold, but defaults, config names, and some behaviors will not. The
text flags the places where it knows the version matters; for everything else,
re-run the relevant experiment on your stack before you rely on a number.

\section*{Exercises}

Exercises at the end of each chapter ask you to predict first and run second:
write down the plan you expect, the number of shuffles, the number of files
scanned, and only then run the query. The learning is in the gap between the
prediction and the result.
